\section{Implementação}

Esta seção descreve a execução efetiva do protótipo e as decisões tomadas durante sua implementação. Diferentemente da seção de Projeto, que apresenta o fluxo planejado, esta seção registra os componentes utilizados, a mudança de rota de integração e os artefatos preservados para reprodução. Os resultados quantitativos, as observações qualitativas e suas interpretações são apresentados exclusivamente na seção de Resultados; esta seção não antecipa conclusões de desempenho ou fidelidade.

\subsection{Arquitetura efetivamente implementada}

A Figura \ref{fig:arquitetura_desenvolvida} resume a arquitetura utilizada no estudo de caso \textit{poster}. Ela separa a preparação e avaliação no Google Colab, os artefatos persistidos, a integração no protótipo Unity e a coleta nativa no Meta Quest 3S. Essa separação também delimita os ambientes de medição: métricas de imagem são produzidas a partir das renderizações offline, enquanto desempenho, contadores disponíveis de memória e observações visuais pertencem à execução no headset.

\begin{figure}[H]
\centering
\caption{Arquitetura efetivamente implementada no estudo de caso \textit{poster}.}
\label{fig:arquitetura_desenvolvida}
\resizebox{0.97\textwidth}{!}{%
\fbox{%
\begin{tikzpicture}[
    bloco/.style={rectangle, rounded corners, draw, align=center, text width=3.15cm, minimum height=1.05cm, inner sep=5pt, font=\footnotesize},
    artefato/.style={rectangle, draw, align=center, text width=3.15cm, minimum height=0.9cm, inner sep=4pt, font=\footnotesize, fill=black!6},
    seta/.style={-{Latex[length=2.2mm]}, thick}
]
\node[bloco] (dadosreal) {Dataset \textit{poster}\\imagens RGB, poses e divisão de teste};
\node[bloco, right=0.8cm of dadosreal] (colabreal) {Google Colab / Nerfstudio\\treinamento, avaliação offline e exportação};
\node[artefato, below=0.85cm of colabreal] (artefatosreal) {Configurações, PLYs, hashes, manifestos e registros};
\node[bloco, right=0.8cm of colabreal] (unityreal) {Protótipo Unity\\importação, XR, Android/Vulkan e APKs};
\node[bloco, below=0.85cm of unityreal] (questreal) {Meta Quest 3S\\pose controlada, métricas, capturas e observação};
\node[artefato, below=0.85cm of artefatosreal] (consolidadoreal) {Consolidação derivada\\CSV, tabelas e gráficos reproduzíveis};
\draw[seta] (dadosreal) -- (colabreal);
\draw[seta] (colabreal) -- (unityreal);
\draw[seta] (colabreal) -- (artefatosreal);
\draw[seta] (artefatosreal) -- (questreal);
\draw[seta] (unityreal) -- (questreal);
\draw[seta] (artefatosreal) -- (consolidadoreal);
\draw[seta] (questreal) -- (consolidadoreal);
\end{tikzpicture}%
}}
\legend{Fonte: Elaborado pelo autor.}
\end{figure}

Os arquivos derivados organizam os dados para consulta e geração de tabelas, mas não substituem os manifestos nem as medições brutas. Por essa razão, uma captura da Unity ou uma medida de GPU do desktop não é reutilizada como métrica offline nem como medição de memória do Quest.

\subsection{Implementação do pipeline experimental}

O estudo de caso adotado foi o conjunto \textit{poster} disponibilizado pelo Nerfstudio. As etapas de pré-processamento, treinamento, avaliação e exportação foram executadas no Google Colab, com os modelos e registros de configuração preservados em armazenamento externo. O modelo NeRF foi mantido como referência de fidelidade visual, enquanto as representações 3D-GS foram treinadas ou derivadas a partir da mesma cena para a avaliação de custo e desempenho.

A execução foi organizada de modo incremental. Primeiro, foi preservada uma \textit{baseline} Splatfacto; depois, foram produzidas variantes de regularização de escala, de maior capacidade e de poda determinística por opacidade. As podas de 100 mil e 50 mil gaussianas foram geradas a partir do PLY congelado da \textit{baseline}, mantendo os registros binários completos das gaussianas selecionadas. Cada artefato possui identificador de variante, contagem de gaussianas, tamanho e soma de verificação registrados em seus manifestos.

\subsection{Poda determinística por opacidade}

As variantes podadas foram geradas por seleção determinística dos registros binários completos do PLY da \textit{baseline}. Os \textit{logits} de opacidade armazenados são ordenados de forma decrescente; empates são resolvidos pelo índice original do registro. Após a seleção, os registros são emitidos na ordem original, preservando posição, coeficientes de cor, opacidade, escala e rotação. A estratégia foi adotada para isolar a redução de gaussianas sem introduzir novo treinamento.

O Algoritmo \ref{lst:poda_opacidade} apresenta a lógica central. O pseudocódigo sintetiza a seleção e a preservação dos parâmetros das gaussianas, aspectos necessários à validade da comparação.

\begin{lstlisting}[caption={Pseudocódigo da poda determinística por opacidade.},label={lst:poda_opacidade}]
procedure GenerateOpacityTopK(sourcePly, k)
    header, count, stride, opacityOffset <- ReadAndValidateHeader(sourcePly)
    ranked <- empty list
    for index <- 0 to count - 1 do
        record <- ReadRecord(sourcePly, stride)
        opacity <- ReadFloat(record, opacityOffset)
        Append(ranked, (opacity, index))
    end for
    selected <- FirstK(SortByOpacityDescendingThenIndexAscending(ranked), k)
    WriteHeaderWithVertexCount(outputPly, k)
    for each record in sourcePly in source order do
        if Index(record) belongs to selected then
            WriteCompleteRecord(outputPly, record)
        end if
    end for
    SaveReport(sourcePly, outputPly, selected, SHA256(sourcePly), SHA256(outputPly))
end procedure
\end{lstlisting}

As podas não foram tratadas como modelos treinados novamente. Na avaliação offline, a configuração da \textit{baseline} foi carregada e seus parâmetros gaussianos foram substituídos pelo conteúdo dos PLYs derivados nas mesmas poses e imagens retidas. Assim, a alteração da representação e a proveniência de cada arquivo permaneceram explícitas.

\subsection{Mudança da rota de integração}

Durante a implementação inicial, foi investigada uma rota de integração com a Unreal Engine. Essa rota não foi consolidada para execução nativa no headset devido aos riscos de compatibilidade entre versões, plugins experimentais, importação de gaussianas e geração de \textit{builds} Android/Vulkan. A alteração para Unity não é interpretada como comparação de superioridade entre motores gráficos; foi uma decisão metodológica para preservar uma rota nativa, estável e reproduzível dentro do escopo do trabalho.

A integração implementada passou a utilizar Unity 6000.5.8f1, Universal Render Pipeline, OpenXR, Android/Vulkan e UnitySplats. O projeto Unity importa os PLYs das variantes como ativos independentes e mantém uma cena mínima de avaliação, com carregamento do modelo, câmera estereoscópica e configuração controlada de pose. As versões do motor, do pacote, do \textit{build} e os identificadores de representação são preservados nos manifestos e nos registros de \textit{build}.

\subsection{Execução no Meta Quest 3S e preservação de evidências}

Os APKs foram gerados para Android/Vulkan e instalados no Meta Quest 3S em execução nativa. A coleta utilizou uma condição estacionária de referência, com alinhamento \textit{FullPoseOnce}: após o rastreamento, a pose completa é fixada para reduzir a variação entre repetições. Essa condição mede o custo de renderização estacionária e não constitui avaliação de locomoção, conforto durante movimento ou comportamento térmico prolongado.

Cada execução retida possui metadados do APK, identificador esperado de variante, relatório de medições da aplicação, marcador de pose e, quando disponível, arquivo CSV do OVR Metrics. Capturas e observações visuais são mantidas como evidência qualitativa separada das métricas offline. A execução por desktop via Meta Horizon Link foi também realizada como contingência diagnóstica em Windows, com \textit{build} Direct3D12 e IL2CPP; seus registros permanecem separados da série Android e são apresentados apenas como resultado complementar.

\subsection{Testes preliminares e contingências}

Os testes preliminares ajudaram a definir a configuração comparativa final, mas não foram agregados à série principal. Foram realizados diagnósticos de modo monoscópico, ordenação por \textit{renderer}, alinhamento de pose e maior capacidade da representação. Nesses testes, registrei artefatos de fundo, cintilação e limitações de desempenho que motivaram a padronização posterior da condição \textit{FullPoseOnce} e das repetições comparáveis. Essas observações orientaram decisões de implementação, mas não foram tratadas como medidas comparáveis da série final.

A execução nativa no Quest permaneceu o cenário-alvo. A renderização em desktop com transmissão para o headset foi mantida como contingência diagnóstica, separada metodologicamente da série Android. A tentativa Linux com WiVRn foi preservada como falha: embora o pareamento do headset tenha sido concluído, o jogador Unity não carregou o componente OpenXR, não recebeu rastreamento do HMD e não estabeleceu uma sessão VR válida. Logo, ela não gerou resultado comparável nem substituiu a avaliação \textit{standalone}.

\subsection{Rastreabilidade e limitações de implementação}

Os manifestos do treinamento, das podas, dos APKs e das execuções permitem relacionar cada resultado a uma cena, uma representação e uma condição de teste. A consolidação derivada não altera os dados brutos e preserva a distinção entre a série comparativa de cinco execuções e diagnósticos isolados. Essa distinção é necessária para que uma observação pontual de maior fidelidade ou uma falha de desempenho não seja apresentada como média comparável.

As podas determinísticas foram avaliadas offline nas mesmas imagens de teste por substituição dos parâmetros gaussianos na configuração da baseline. Essa avaliação foi mantida separada das capturas da Unity e do Quest. Persiste a limitação de que a condição de pose fixa não caracteriza estabilidade durante movimento.

Para verificar a execução integral do procedimento, executei novamente todas as etapas em um ambiente Google Colab limpo, utilizando uma GPU Tesla T4 e 30 mil passos por modelo. Concluí as sete etapas de treinamento, avaliação e exportação na primeira tentativa e conferi o tamanho e a soma de verificação SHA-256 de nove artefatos produzidos. Nessa validação, reutilizei as poses existentes do conjunto \textit{poster}, sem executar uma nova estimativa por COLMAP. Preservei os modelos resultantes como evidência de validação operacional e mantive os modelos originais na comparação experimental.

A descrição dos procedimentos próprios concentra-se na arquitetura, nas decisões técnicas e nas consequências metodológicas da implementação. O pseudocódigo apresentado sintetiza a operação central da poda e permite compreender como os parâmetros das gaussianas selecionadas foram preservados.
